\slajdid{02-s004}
\begin{frame}[label=02-s004]
\gusttekst
\centering
Zapamtiti\qquad\fcolorbox{DEcrvena}{white}{\(\displaystyle t_{pHL}\ne t_{pLH}\)}

Uvešćemo pojam i srednjeg kašnjenja = kašnjenje logičkog kola
\[
t_p=\frac{t_{pHL}+t_{pLH}}{2}
\]
i radi jednostavnosti posmatranja \alert{nekih} pojava smatrati
\[
t_{pHL}=t_{pLH}=t_p
\]
A onda često pogrešno razmišljanje za \alert{sve} slučajeve

\begin{columns}[c,onlytextwidth]
\begin{column}{.46\textwidth}\centering
\begin{tikzpicture}[DEoznaka,x=10mm,y=7.5mm]
\draw[DEstrelica] (-.7,1.9)--(3.8,1.9) node[below] {\(t\)};
\draw[DEstrelica] (0,1.75)--(0,3.7) node[left] {\(x(t)\)};
\draw[DEstrelica] (-.7,-.1)--(3.8,-.1) node[below] {\(t\)};
\draw[DEstrelica] (0,-.25)--(0,1.7) node[left] {\(y(t)\)};
\node[above left,inner sep=1pt] at (0,3) {=1};\node[above left,inner sep=1pt] at (0,2.1) {=0};
\node[above left,inner sep=1pt] at (0,1) {=1};\node[above left,inner sep=1pt] at (0,.1) {=0};
\draw[DEvod] (-.5,2.1)--(.6,2.1)--(.6,3)--(.85,3)--(.85,2.1)--(3.5,2.1);
\draw[DEvod] (-.5,1)--(2.8,1)--(2.8,.1)--(3.05,.1)--(3.05,1)--(3.5,1);
\draw[dashed] (.6,3)--(.6,-.35) (2.8,.1)--(2.8,-.35) (0,.1)--(2.8,.1);
\draw[<->] (.6,-.3)--(2.8,-.3) node[midway,below] {\(t_p\)};
\end{tikzpicture}

\end{column}
\begin{column}{.52\textwidth}\gusttekst
Signal na izlazu će biti isti samo zakašnjen?

Dolazi do promene izlaznog signala, trajanja,\\
a moguće da se neće ni pojaviti!

\centering
Zavisi od\\
Trajanje signala PREMA kašnjenjima
\end{column}
\end{columns}
\end{frame}
